\section{Arithmetic reduction and conventions}\label{sec:reduction}

The converse at $2$ settles the even prime.  For an odd prime, we
choose auxiliary quadratic fields so that the problem reduces to
proving that one conductor-one Heegner trace is non-torsion.  The
resulting compact Selmer spaces will also supply the initial data
for the tame deformation.

\subsection{Selmer groups and established inputs}

For a number field $M$, put $T_pA_0=\varprojlim_n A_0[p^n]$ and
$V_pA_0=T_pA_0\otimes_{\Z_p}\Q_p$. For a finite place $v$ of $M$,
the finite local condition is the Kummer image
\[
 \Hf(M_v,V_pA_0)
 =\im\bigl(A_0(M_v)^{\wedge}_p\otimes_{\Z_p}\Q_p
       \longrightarrow H^1(M_v,V_pA_0)\bigr).
\]
Here $A_0(M_v)^{\wedge}_p=\varprojlim_n A_0(M_v)/p^nA_0(M_v)$.
The rational compact Selmer group is
\[
 H^1_f(M,V_pA_0)=
 \ker\left(H^1(M,V_pA_0)\longrightarrow
   \prod_v H^1(M_v,V_pA_0)/\Hf(M_v,V_pA_0)\right).
\]
Continuous cohomology and the usual archimedean local conditions are
understood. The Kummer exact sequences and local/global duality give
\begin{equation}\label{eq:compact-corank}
 \dim_{\Q_p}H^1_f(M,V_pA_0)
 =\corank_{\Z_p}\Sel_{p^\infty}(A_0/M).
\end{equation}
We use rational coefficients throughout the cohomological argument;
integral lattices will be retained whenever a uniform valuation bound
is required.

We use the following two established results.

\begin{theorem}[Established converse at $2$]
\label{thm:two-converse-input}
If $E/\Q$ is any elliptic curve and
$\corank_{\Z_2}\Sel_{2^\infty}(E/\Q)=r\in\{0,1\}$, then
\[
 \ord_{s=1}L(E,s)=\rank E(\Q)=r,
 \qquad \Sha(E/\Q)\text{ is finite}.
\]
\end{theorem}

\begin{theorem}[Established analytic twist-density theorem]
\label{thm:twist-density-input}
For any elliptic curve $E/\Q$, among nonzero squarefree integers $d$
ordered by $|d|$, the quadratic twists $E^{(d)}$ of analytic rank zero
and one each have density $1/2$. In particular the exceptional set of
twists of analytic rank at least two has density zero.
\end{theorem}

These are \citet[Theorems~1.1 and~1.2]{OpenAIGoldfeldDensity2026}.
Theorem~\ref{thm:two-converse-input} treats the prime $2$.
For auxiliary-field selection we use the density-zero exceptional-set
consequence of Theorem~\ref{thm:twist-density-input} in fixed
positive-density congruence classes.
The remainder of the paper supplies the argument at odd primes.

Fix $A_0$ and $r=s_p(A_0)\in\{0,1\}$. By
Theorem~\ref{thm:two-converse-input}, we may assume $p>2$. Write
$N$ for the conductor and $V=V_pA_0$. The Weil pairing gives
$V\simeq V^*(1)$ with alternating pairing. The $p$-parity theorem
\citep[Theorem~1.4]{DokchitserDokchitser2010} gives the global root
number
\begin{equation}\label{eq:root-sign}
 w(A_0)=(-1)^r.
\end{equation}
Modularity and the Gross--Zagier--Kolyvagin theorems will be used in
their rank-zero and rank-one forms
\citep{BCDT2001,GrossZagier1986,Kolyvagin1990}.

\subsection{Two auxiliary discriminants}

\begin{lemma}\label{lem:auxiliary-fields}
There are coprime negative odd fundamental discriminants $D,D'$ such
that every prime dividing $2Np$ splits in both quadratic fields,
\[
 \left(\frac{D}{|D'|}\right)=1,
\]
and each of $A_0^D,A_0^{D'},A_0^{DD'}$ has the smallest analytic rank
permitted by its sign. They may also be chosen so that
\[
 K=\Q(\sqrt D),\qquad F=\Q(\sqrt{DD'}),\qquad
 L=\Q(\sqrt D,\sqrt{D'})
\]
satisfy the following properties: $L/F$ is CM and unramified at every
finite prime, $L$ has only the roots of unity $\{\pm1\}$, and $L$
does not contain the CM field of $A_0$ if $A_0$ has CM.
\end{lemma}

\begin{proof}
Negative odd fundamental discriminants in any admissible fixed
congruence class have positive density. The splitting conditions at
$2Np$ are finitely many such conditions. We can exclude the finitely
many unwanted quadratic fields and then apply
Theorem~\ref{thm:twist-density-input} to choose $D$ for which $A_0^D$
has analytic rank at most one. Its rank is then minimal for its sign.

Keep $D$ fixed. The conditions on $D'$ of being coprime to $D$,
splitting $2Np$, and satisfying the displayed Jacobi-symbol condition
again contain admissible congruence classes of positive density. For
example, impose $D'\equiv-1\pmod{|D|}$ and $D'\equiv1$ at a
sufficiently divisible modulus supported on $2Np$. The moduli are
coprime, and $|D'|\equiv1\pmod{|D|}$ gives the required Jacobi
symbol. Exclude the finitely many remaining forbidden fields. Apply
the density theorem to both $A_0$ and $A_0^D$. The union of their two
exceptional sets has density zero, so some $D'$ in the prescribed
classes makes both $A_0^{D'}$ and $(A_0^D)^{D'}=A_0^{DD'}$ have
analytic rank at most one.

This use of density is unchanged if parameters are described by
fundamental discriminants: for squarefree $d$ the corresponding
discriminant is $d$ or $4d$, so a zero-density exceptional set remains
zero-density. In the odd classes used here the parameter itself is
the discriminant.

The three quadratic subfields have discriminants $D,D',DD'$.
Consequently the biquadratic discriminant formula gives
\[
 |\operatorname{disc}(L)|=D^2(D')^2
   =\operatorname{disc}(F)^2.
\]
The discriminant formula for $L/F$ therefore has relative
discriminant of norm one, proving finite unramifiedness. The two
imaginary subfields give the CM involution over the real field $F$.
The exclusions made above remove any extra roots of unity and the
possible CM field of $A_0$.
\end{proof}

\begin{proposition}\label{prop:arithmetic-reduction}
With these choices,
\[
 \dim H^1_f(K,V)=1,
 \qquad
 \dim H^1_f(K,V^{DD'})=1.
\]
The second line localizes nontrivially at both places of $K$ over $p$.
Let $y\in A_0(K)$ be the norm of the conductor-one Heegner point
under a modular parametrization taking the cusp $\infty$ to zero.
If $y$ is non-torsion, Theorem~\ref{thm:main} follows.
\end{proposition}

\begin{proof}
For a fundamental discriminant prime to $N$, the quadratic-twist
root-number formula is $w(A_0^d)=w(A_0)\chi_d(-N)$.
The splitting conditions and the signs of $D,D'$ therefore give
\[
 w(A_0^D)=w(A_0^{D'})=-(-1)^r,
 \qquad w(A_0^{DD'})=(-1)^r.
\]
By Lemma~\ref{lem:auxiliary-fields}, their analytic ranks are
$1-r,1-r,r$, respectively. Gross--Zagier--Kolyvagin gives the same
Mordell--Weil and $p$-power Selmer ranks and finite Sha for these three
twists. Restriction and corestriction for the quadratic extension
$K/\Q$ give
\[
 H^1_f(K,V)
  =H^1_f(\Q,V)\oplus H^1_f(\Q,V^D),
\]
and the analogous decomposition for $V^{DD'}$. Thus the two dimensions
are $r+(1-r)=1$ and $r+(1-r)=1$.

For the second representation both summands come from twists with
known analytic rank at most one. Its Selmer line is therefore spanned
by a non-torsion global point on the corresponding twist over $K$.
Such a point has nonzero image under the local logarithm at either
place above $p$. Indeed, on a fixed $p$-adic completion, a finite-index
open subgroup is a formal group on which the logarithm is injective
after shrinking; a point with zero logarithm consequently has a
torsion multiple in that subgroup, and is torsion. This argument uses
no reduction hypothesis.

Every prime dividing $N$ splits in $K$, so the classical Heegner
hypothesis holds. The Gross--Zagier formula identifies the nonzero
height of $y$ with simple vanishing of
\[
 L(A_0/K,s)=L(A_0,s)L(A_0^D,s)
\]
at $s=1$. Since the second factor has order $1-r$, a non-torsion $y$
implies $a(A_0)=r$. The analytic rank-zero/one theorems then give
$\rank A_0(\Q)=r$ and finiteness of the whole $\Sha(A_0/\Q)$.
\end{proof}

For the rest of the proof we assume
\begin{equation}\label{eq:torsion-assumption}
 y=0\quad\text{in }A_0(K)\otimes\Q.
\end{equation}
We will derive a contradiction. In particular, every subsequent
choice is made after $A_0,p,D,D'$ have been fixed.
The first Selmer line in Proposition~\ref{prop:arithmetic-reduction}
may have zero localization, whereas the second cannot.
Section~\ref{sec:tame} distinguishes these two possibilities before
deforming the representations.

\subsection{Coefficient, character, and measure conventions}
\label{sec:conventions}

Fix embeddings into $\C$ and $\overline{\Q}_p$. Let $\Sigma$ be the CM
type on $L$ induced from one embedding of $K$, and let $\Sigma_p$ be
the corresponding half of the places of $L$ above $p$. Complex
conjugation is denoted by $c$. All the fields introduced above are
split at $p$, and the places of $K$ above $p$ will be distinguished
according to this choice.

Smooth characters and normalized Satake parameters are evaluated on
uniformizers using the reciprocity convention corresponding to
geometric Frobenius. Write $\epsilon$ for the cyclotomic character.
A de Rham weight of $\epsilon^{-1}$ is called $1$. When a local
cohomology computation instead uses arithmetic Frobenius, this will
be stated; inversion changes neither the vanishing assertions nor
the invertibility requirements used below.

Characters of $\U(1)_{L/F}=\ker\Norm_{L/F}$ are pulled back to
characters of $L^\times\backslash\A_L^\times$ by $z\mapsto z/z^c$;
the same convention applies to $K/\Q$. Choose a character
$\mathcal C$ of $\U(1)_{K/\Q}$ with infinity type $s\mapsto s$ in
the chosen embedding and with finite conductor supported on split
primes away from $pDD'N$. These primes may be required to split
completely in $L$. The ray-class construction gives such a character:
a sufficiently deep split conductor removes the finite unit
obstruction, and the prescribed algebraic character on principal
ideles then extends across the finite ray-class quotient. We use the
same symbol for its norm pullback to $\U(1)_{L/F}$ and for the
corresponding algebraic $p$-adic character.

At the distinguished split component above $p$, its finite smooth
part has uniformizer value of valuation $-1$ with these conventions.
This records the algebraic infinity contribution; it is not a claim
that the smooth character is unitary in the $p$-adic norm. All
archimedean weights $m\to+\infty$ below also satisfy $m\to0$
$p$-adically on a fixed residue branch, including the finite torsion
orders of the characters raised to the $m$th power.

Let $k$ be a sufficiently large fixed finite extension of $\Q_p$,
with ring of integers $\cO$. We enlarge it finitely when fixed
algebraic normalization factors require this. Valuation bounds refer
to fixed bases and integral lattices. A \emph{uniform} bound may
depend on the fixed data and on a fixed truncation degree, but not on
the moving prime, its growing $p$-power quotient, or the growing
weights. When coefficient fields for classical eigensystems vary,
we retain their embeddings in $\overline{\Q}_p$ and use valuation
bounds; no compactness of their union is assumed.

For an algebraic group $H$ over a number field $M$, write
$[H]=H(M)\backslash H(\A_M)$.

Measures on compact source groups and archimedean integration tori
have algebraic normalization, and maximal compact subgroups have
volume one at almost every finite place. Passing between these
measures and invariant norm formulas costs a fixed global scalar.
All toric sums called \emph{raw} below are sums without division by
the varying ring-class order. This distinction is essential for
the estimates at the moving primes.
